\section{A complete five-loop source-current family}
\label{currents:family}

\subsection{Original graph and retained source support}

The example has seven vertices, labelled $0,1,\ldots,6$, and the following
ordered, oriented edges:
\begin{equation}
 (01,13,02,24,05,06,35,36,45,46,56).
 \label{currents:edges}
\end{equation}
The two marked vertices are $5,6$. The independent source spinors are
$\lambda_1,\lambda_2,\lambda_5$, with the remaining sources fixed by
\begin{equation}
 \lambda_0=\lambda_3=\lambda_4=0,\qquad
 \lambda_6=-\lambda_1-\lambda_2-\lambda_5.
 \label{currents:sources}
\end{equation}
The quartic vertex is $0$; its zero source is part of this family.
Sources at vertices $0,3,4$ are not varied by the result.

Write $K(\lambda)$ for the normalized integrated bare kernel and
$M_{e,+\dot a}(\lambda)$ for the integral with the complete original
current $y_{e,+\dot a}$ inserted on edge $e$. The index $+\dot a$
remains free. With the incidence sign equal to $+1$ at an edge's
initial vertex and $-1$ at its terminal vertex, current conservation
requires, before and after integration,
\begin{equation}
 \sum_{e:\,e(0)=v}M_{e,+\dot a}
 -\sum_{e:\,e(1)=v}M_{e,+\dot a}
 =\lambda_{v,+\dot a}K.
 \label{currents:conservation}
\end{equation}
All seven spinor equations are kept, including the dependent total sum.

\begin{theorem}[Complete supported-source polynomial]
\label{currents:complete}
For the ordered graph \eqref{currents:edges} and source support
\eqref{currents:sources}, the bare kernel and all eleven first-current
spinors are homogeneous source polynomials of degrees ten and eleven.
Their coefficients belong to $\Q+\Q\zeta(3)+\Q\zeta(5)$. The complete
formulas have sixteen nonzero bare bracket monomials and 468 nonzero
current bracket/source monomials, with 1,334 nonzero rational
coefficients in the three displayed period channels. Every bracket
power is nonnegative, so the formulas include all degenerate source
frames. The complete coefficient tables are printed in the companion
supplement and recorded in the exact certificate.
\end{theorem}

The proof has two parts: a finite algebraic reconstruction whose
completeness is proved a priori, and exact integration of every sample
needed by that reconstruction. We describe both, including the
independent controls, rather than infer a general formula from a
numerical fit.

\subsection{Polynomial degree and coefficientwise convergence}

The five-loop current-area form contains ten source-linear current
differentials. Its complete rational numerator therefore has source
degree ten, and its denominator is independent of all sources. Inserting
one full original current raises the degree by one. All tree-flow
currents are finite source sums; the full network current is a convex
combination of them at positive resistances. Thus each inserted current
is bounded on positive projective resistance space at a fixed source
assignment.

The compact-cube construction of the preceding network section proves
absolute convergence of the complete density at each fixed rational
balanced source assignment. Multiplication by a bounded current
preserves it. To justify integrating every source coefficient, take a
finite rational unisolvent grid for polynomials of total degree at most
eleven in the six independent source coordinates. Each coefficient
integrand is a finite rational linear combination of evaluated
integrands on that grid. The triangle inequality proves absolute
convergence of each coefficient. Integration is consequently a
well-defined linear map on the complete finite source polynomial.

For an invertible linear change $g$ of the two-dimensional source
space, the alternating spinor pairing is multiplied by $\det g$.
The five area factors and the remaining free current index give
\begin{align}
 K(g\lambda)&=(\det g)^5 K(\lambda),\label{currents:scalarcov}\\
 M_{e,+\dot a}(g\lambda)&=(\det g)^5
 g_{+\dot a}{}^{+\dot b}M_{e,+\dot b}(\lambda).
 \label{currents:vectorcov}
\end{align}
These identities hold before integration. They are preserved by all
rational chain integrations and scalar-period reductions. No
independence assumption about the numerical constants is involved.

\subsection{An exact interpolation bound}

Choose a source frame with $\lambda_1=(1,0)$,
$\lambda_2=(0,1)$ and $\lambda_5=(a,b)$. The scalar coefficient
functions of a current in this frame are polynomials $f_e(a,b)$,
$g_e(a,b)$. We now prove the bound on their degrees, rather than
assume one from observed samples.

Temporarily set $\lambda_1=(1,0)$, $\lambda_2=(s,t)$,
$\lambda_5=(0,1)$. For $t\ne0$, source covariance gives the two
frame coefficient functions
\begin{equation}
 t^5\bigl(f_e(-s/t,1/t)+s\,g_e(-s/t,1/t)\bigr),\qquad
 t^6g_e(-s/t,1/t).
 \label{currents:degreeproof}
\end{equation}
The original source expression is polynomial in $s,t$, including
$t=0$. If a homogeneous part had degree greater than six, the second
expression would force that highest part of $g_e$ to vanish, and then
the first would force the corresponding part of $f_e$ to vanish.
Descending on degree proves
\begin{equation}
 \deg f_e\le6,\qquad \deg g_e\le6.
 \label{currents:degree}
\end{equation}
At degree six, regularity of the first expression instead gives
\begin{equation}
 b f_{e,6}(a,b)=a g_{e,6}(a,b).
 \label{currents:syzygy}
\end{equation}
Here the subscript $6$ denotes the homogeneous part of degree six.
Because the two coordinate polynomials $a,b$ are relatively prime,
there is a homogeneous polynomial $h_{e,5}$ of degree five with
\begin{equation}
 f_{e,6}=a h_{e,5},\qquad g_{e,6}=b h_{e,5}.
 \label{currents:high}
\end{equation}
The scalar version of \eqref{currents:degreeproof}, with only the
factor $t^5$, proves degree at most five for the bare frame polynomial.

There are $\binom{8}{2}=28$ monomials $a^i b^j$ with $i+j\le6$.
The 28 integer nodes
\begin{equation}
 a,b\ge0,\qquad a+b\le6
 \label{currents:nodes}
\end{equation}
are unisolvent: the basis of products of falling factorials has a
triangular evaluation matrix with diagonal $i!j!$. The production
reconstruction inverts the monomial evaluation matrix over $\Q$.
The independent verifier instead uses
\begin{equation}
 \frac{\Delta_a^i\Delta_b^j F(0,0)}{i!j!}
 =\frac1{i!j!}\sum_{r=0}^{i}\sum_{s=0}^{j}
 (-1)^{i+j-r-s}\binom{i}{r}\binom{j}{s}F(r,s),
 \label{currents:newton}
\end{equation}
and expands the falling-factorial basis. The differences in this
formula are forward differences of the full scalar coefficient
polynomial. Every node on its right satisfies $r+s\le i+j\le6$.

The graph involution $(1\ 2)(3\ 4)$ induces the even edge permutation
$(0\ 2)(1\ 3)(6\ 8)(7\ 9)$ in the zero-based order
\eqref{currents:edges}; no edge orientation reverses. Combining it
with the determinant-minus-one exchange of the source coordinates
takes $(a,b)$ to $(b,a)$. It changes the bare sign and acts on the
full current by the corresponding source-coordinate exchange and
minus sign. Hence the sixteen nodes $a\le b$ supply the remaining
twelve nodes exactly. The symmetry is a derived equality of original
graphs, not an assumption about an integrated period table.

\subsection{Exact scalar integration of the samples}

The sourced chains are $0$--$1$--$3$ and $0$--$2$--$4$.
Their exact simplex integrations leave the complete graph on
vertices $0,3,4,5,6$ with the edge $34$ absent. Each scalar monomial
retains its exact factorial coefficient and original-current label.
Projective inversion maps this scalar family to the planar-dual
triangular prism, replacing an edge power by the denominator power
minus that edge power. The existing exact prism tables at denominator
powers seven and eight evaluate every bare and first-current
monomial encountered here. Strict convergence is checked before each
lookup; an absent table entry is an error, never a zero coefficient.

At every direct source node, five complete chord-current spinors are
integrated independently. The six tree currents follow from the
original tree flow plus their signed fundamental-cycle contributions.
In particular the tree-flow term is multiplied by the full bare
kernel and retained. This reconstructs all eleven full spinors and
gives \eqref{currents:conservation} at every node. There are
$28(1+22)=644$ frame coefficient dictionaries. Both reconstruction
methods agree on every dictionary and on every coefficient of all
fourteen scalar entries of \eqref{currents:conservation}.

\subsection{The denominator-free covariant formula}

For two original spinors, $\langle\lambda_u,\lambda_v\rangle$
denotes their alternating determinant. On the open set where
$\langle\lambda_1,\lambda_2\rangle\ne0$, the frame coordinates are
\begin{equation}
 a=\frac{\langle\lambda_5,\lambda_2\rangle}
 {\langle\lambda_1,\lambda_2\rangle},\qquad
 b=\frac{\langle\lambda_1,\lambda_5\rangle}
 {\langle\lambda_1,\lambda_2\rangle}.
\end{equation}
Let $f_{e,ij},g_{e,ij}$ be the coefficients of $a^i b^j$ in the
parts of $f_e,g_e$ of degrees at most five, and let $h_{e,ij}$ be the
coefficients of $h_{e,5}$. Then the complete original free-index result
is
\begin{align}
 M_{e,+\dot a}(\lambda)
 ={}&\sum_{i+j\le5}
 \langle\lambda_1,\lambda_2\rangle^{5-i-j}
 \langle\lambda_5,\lambda_2\rangle^i
 \langle\lambda_1,\lambda_5\rangle^j
 \bigl(f_{e,ij}\lambda_{1,+\dot a}
       +g_{e,ij}\lambda_{2,+\dot a}\bigr)
 \nonumber\\
 &+\sum_{i+j=5}
 \langle\lambda_5,\lambda_2\rangle^i
 \langle\lambda_1,\lambda_5\rangle^j
 h_{e,ij}\lambda_{5,+\dot a}.
 \label{currents:covariant}
\end{align}
The second line follows from
$a\lambda_1+b\lambda_2=\lambda_5$ and
\eqref{currents:high}. Every exponent is nonnegative. The bare formula
uses the first line's bracket monomials without the extra source
spinor. Both sides are polynomial in all original source coordinates
and agree on a dense open set. Thus they agree everywhere, including
parallel or zero root sources. This proves the reconstruction and
degenerate-extension parts of Theorem~\ref{currents:complete}.

Every coefficient in \eqref{currents:covariant} is evaluated by the
finite exact scalar integrations above. The companion supplement
prints the full list; its table generator and independent parser
verify all 468 nonzero current monomials and every one of their
1,334 nonzero rational period coefficients.

\subsection{A complete vanishing-bare control}

Set $\lambda_5=0$ and retain arbitrary $\lambda_1,\lambda_2$.
Every term of the complete bare polynomial contains a bracket with
$\lambda_5$, so $K=0$. Define the two scalar coefficients in each
row of the following table by
\begin{equation}
 M_{e,+\dot a}\big|_{\lambda_5=0}
 =\langle\lambda_1,\lambda_2\rangle^5
 \bigl(A_e\lambda_{1,+\dot a}+B_e\lambda_{2,+\dot a}\bigr).
 \label{currents:barezero}
\end{equation}
The table gives all original edges, with their original orientations.
\begin{center}\small
\begin{tabular}{@{}crr@{}}
\toprule Original edge & $A_e$ & $B_e$ \\ \midrule
$0\to 1$ & $\frac{923}{1828915200}-\frac{1}{8467200}\zeta(3)$ & $\frac{17}{114307200}-\frac{1}{8467200}\zeta(3)$ \\
$1\to 3$ & $\frac{923}{1828915200}-\frac{1}{8467200}\zeta(3)$ & $\frac{17}{114307200}-\frac{1}{8467200}\zeta(3)$ \\
$0\to 2$ & $-\frac{17}{114307200}+\frac{1}{8467200}\zeta(3)$ & $-\frac{923}{1828915200}+\frac{1}{8467200}\zeta(3)$ \\
$2\to 4$ & $-\frac{17}{114307200}+\frac{1}{8467200}\zeta(3)$ & $-\frac{923}{1828915200}+\frac{1}{8467200}\zeta(3)$ \\
$0\to 5$ & $\frac{38987}{30481920000}-\frac{137}{105840000}\zeta(3)$ & $-\frac{38987}{30481920000}+\frac{137}{105840000}\zeta(3)$ \\
$0\to 6$ & $-\frac{49837}{30481920000}+\frac{137}{105840000}\zeta(3)$ & $\frac{49837}{30481920000}-\frac{137}{105840000}\zeta(3)$ \\
$3\to 5$ & $\frac{1597}{2612736000}-\frac{1}{6048000}\zeta(3)$ & $\frac{341}{338688000}-\frac{11}{14112000}\zeta(3)$ \\
$3\to 6$ & $-\frac{1949}{18289152000}+\frac{1}{21168000}\zeta(3)$ & $-\frac{1121}{1306368000}+\frac{1}{1512000}\zeta(3)$ \\
$4\to 5$ & $-\frac{341}{338688000}+\frac{11}{14112000}\zeta(3)$ & $-\frac{1597}{2612736000}+\frac{1}{6048000}\zeta(3)$ \\
$4\to 6$ & $\frac{1121}{1306368000}-\frac{1}{1512000}\zeta(3)$ & $\frac{1949}{18289152000}-\frac{1}{21168000}\zeta(3)$ \\
$5\to 6$ & $\frac{40393}{45722880000}-\frac{1}{1470000}\zeta(3)$ & $-\frac{40393}{45722880000}+\frac{1}{1470000}\zeta(3)$ \\
\bottomrule\end{tabular}\end{center}


For edge $01$, the elementary estimate
\begin{equation}
 \zeta(3)=1+\frac18+\sum_{m=3}^{\infty}\frac1{m^3}
 <1+\frac18+\int_2^{\infty}\frac{dx}{x^3}=\frac54
\end{equation}
gives the exact positive lower bounds
\begin{align}
 \frac{923}{1828915200}-\frac{\zeta(3)}{8467200}
 &>\frac{653}{1828915200},\nonumber\\
 \frac{17}{114307200}-\frac{\zeta(3)}{8467200}
 &>\frac{2}{1828915200}.
\end{align}
Consequently the complete current is a nonzero covariant polynomial
even though the bare polynomial vanishes identically on this
specialization. The complete coefficient table also satisfies every
zero-source divergence equation \eqref{currents:conservation}.

\subsection{Fresh-source and degenerate-frame controls}

Two additional calculations integrate every original edge separately,
using the alternative source-flow tree with edge indices
$(0,2,4,6,8,10)$ instead of the interpolation tree.
Their independent source assignments are
\begin{align}
 (\lambda_1,\lambda_2,\lambda_5)&=((2,1),(1,2),(-1,4)),
 \label{currents:fresh}\\
 (\lambda_1,\lambda_2,\lambda_5)&=((1,0),(2,0),(0,1)).
 \label{currents:degenerate}
\end{align}
The first root determinant is three and its frame point is $(-2,3)$,
outside \eqref{currents:nodes}. The second root determinant is zero.
Each calculation reproduces the bare value and all 22 scalar entries
of the eleven complete spinors. Neither control uses source
interpolation nor chord reconstruction. Their original source
numerators also agree coefficientwise across the two source-flow
and cycle bases. The numerator formula and scalar tables are shared
with the production calculation; the source assignments, network
bases and all-edge integrations are independent checks.

A previous physical-quartic calculation supplies the further source
assignment $\lambda_1=(1,2)$, $\lambda_2=(3,-1)$,
$\lambda_5=(5,3)$ and $\lambda_6=(-9,-4)$. Its complete bare value is
\begin{equation}
 K=\frac{3327759589}{46656000}
   +\frac{50942017}{324000}\zeta(3)
   -\frac{45619}{180}\zeta(5).
 \label{currents:physicalvalue}
\end{equation}
The full reconstructed polynomial gives exactly the same three
rational coefficients.

The different five-loop source family whose suppressed core is
$K_{3,3}$ with one additional edge requires a missing denominator-power-eight
current family. The present theorem concerns the graph
\eqref{currents:edges}; it supplies no result for that other core or
for a complete interacting Wick/colour sum.
